In this study we examine whether the language a chief executive uses in the unscripted question-and-answer session of an earnings call tracks the disclosure state of an acquisition the firm is still withholding, using a residual measure of chief-executive question-and-answer uncertainty that nets out each executive's persistent speaking style.

We find that residual uncertainty is elevated in the quarter before a cash acquisition, with no comparable rise before stock acquisitions; that it becomes indistinguishable from zero once the deal is announced, even before it closes, while the cash that funds it persists to completion; and that the elevation concentrates in cash acquisitions rather than stock on a formal pooled test.

The pattern is not accounted for by analysts devoting more of the call to cash questions; the pre-announcement elevation survives a direct measure of analyst cash scrutiny and its interaction with the pre-announcement window. A further analysis finds the residual unrelated to the firm's post-call bid-ask spread, while the scripted presentation is contemporaneously positively associated with it, consistent with the scripted and unscripted segments relating to the outside information environment differently. The differential-timing result is consistent with treating a withdrawal as a resolution of the deal, and it holds when the cash specification is estimated without its dynamic term; and the run-up and cash-concentration hold when every deal a firm makes is stacked.

Taken together, these patterns suggest that the unscripted language of earnings calls carries a readable, anticipatory trace of a deal's passage from private to public, a signature researchers can in principle observe before the announcement. For the academic literature, the pattern connects two strands usually studied apart: the theory of when an informed manager withholds private information, and the evidence that markets anticipate corporate transactions before they are announced. The within-firm language signature we characterize sits in the window between them. We draw no prescription for investors, managers, or regulators from a correlational pattern that identifies no causal channel; the contribution is to characterize a regularity, not to recommend acting on it.

These findings carry clear limitations. The evidence is correlational and within-firm, and the design supports no causal identification and establishes no mechanism. The concentration in cash deals is motivated in Section~2 but not identified, and the war-chest channel behind cash accumulation in particular remains unestablished. Separately, the source of the uncertainty, a compliance-constrained inability to speak versus a strategically chosen reticence, remains observationally equivalent.

Two further boundaries qualify the evidence. The residual is estimated only for chief executives with enough calls to fix a speaking style, which skews the language sample toward larger, more heavily covered firms, and the scrutiny rule-out is a failure to find rather than a powered test of equivalence, since most calls draw no cash scrutiny at all. The measure itself bounds the inference further. Uncertainty is captured by applying a finance-specific word list to the transcripts, a count of words that abstracts from their context, and the residual is one operationalization of call-specific uncertainty rather than a direct reading of what a chief executive knows. The stock-financed comparison shares the disclosure bind but is an imperfect counterfactual, and the sample, United States public firms from 2002 to 2018, approximately the S\&P~1500, may not extend to other periods, other markets, or the smaller firms outside it.

Future work could ask whether the same anticipatory language signature appears around other classes of withheld material events, and could pursue designs able to separate compliance-constrained from strategic silence, which our setting cannot distinguish. Several extensions follow. The same residual-uncertainty reading could be applied to other corporate transactions, such as divestitures, joint ventures, and strategic alliances, and to settings outside the United States; richer measures of spoken uncertainty could replace the word-list count; and identifying the mechanism behind the cash concentration, which Section~2 motivates but does not establish, including the cash-accumulation channel this study leaves open, would require a design the present setting does not provide.
